\section{Other errors and the limits of causal attribution}
The explicitly derived models cover resistance mismatch, constant coherent
frame rotation, additive voltage/current/speed perturbations, current relabeling,
and a simple timing extension. Gain errors can be expressed locally as
$\vect\varepsilon_u=\matr G_u\voltage+\vect b_u$ and
$\vect\varepsilon_i=\matr G_i\current+\vect b_i$, but their parity depends
on those matrices and on the frame in which gains and offsets originate.
These expressions parameterize errors; they do not establish a complete
sensor-identification procedure.

Deadtime and semiconductor drops contribute structured terminal-voltage errors
that can depend on phase-current signs, switching conditions, and temperature.
Their dq averages are not generally constant sensor offsets. Work on magnetic
map identification explicitly addresses inverter nonlinearities and uncertain
resistance \cite{liu2018}; here they remain nuisance sources unless an
independent voltage model is available. Commanded inverter voltage and actual
terminal voltage must not be silently equated.

\begin{table}[tb]
\centering\small
\begin{tabularx}{\textwidth}{@{}lXXX@{}}
\toprule
Effect & Forbidden d channel & Forbidden q channel & Speed behavior\\
\midrule
Resistance mismatch & Odd, $-i_q\resistanceerror/\omegae$ & Even, $i_d\resistanceerror/\omegae$;
zero at $i_d=0$ & $1/\omegae$ for fixed mismatch\\
Constant angle & Odd angle sensitivity & Even angle sensitivity & No explicit
speed factor in quasi-static map\\
Saturation alone & None under reflection & None under reflection & Magnetic
map held fixed\\
Saturation with angle & Changes odd sensitivity & Changes even sensitivity &
Through operating point or changing magnetics\\
Paired scalar speed error & None for ideal parity & None for ideal parity &
Common multiplicative flux error\\
Voltage/inverter error & Potentially odd & Potentially even & Voltage error
divided by speed; voltage error itself may vary\\
\bottomrule
\end{tabularx}
\caption{Conditional fingerprints derived earlier. Zero forbidden residual
does not imply zero reconstruction error. Pair speed and thermal state must
match; channel amplitudes may vanish at special operating points.}
\label{tab:fingerprints}
\end{table}

The remaining limitations are qualitative. Real construction or magnet
asymmetry can break reflection physically. Harmonics and skew require attention
to rotor-position dependence and the averaging convention; symmetric skew
need not destroy effective reflection. Hysteresis makes flux history-dependent.
Iron loss and other dynamic effects can make a speed-independent static map
inadequate even in an apparently steady fundamental operating condition.
Temperature changes both winding resistance and magnetics. Frequency-dependent
effective drops need not equal a constant winding resistance. These effects
cannot all be resolved from two parity channels.

The method also requires measured support on both sides of the reflection.
Interpolation cannot supply missing experimental information, and extrapolated
partners should not be treated as independent observations. Differentiating
noisy data can corrupt angle sensitivities, so a constrained reference and
uncertainty propagation are prerequisites to quantitative use. Correlated
sensor errors and nuisance terms can invalidate nominal least-squares intervals.

\paragraph{Evidence boundary.}
The equations are analytical under their stated assumptions, and the numerical
checks validate their implementation on a known synthetic model. No experimental
accuracy, temperature range, or inverter rejection performance is claimed.
The cited literature supports the model background and related identification
problems; the specific residual derivations are developed here. A dedicated
experimental reference on calibrating an unknown constant encoder offset from
paired saturated maps, and a measured comparison of effective versus DC winding
resistance, remain open source and validation work. No bibliographic placeholder
is used for these gaps.
